%!TeX root=csp-proof.tex
\section{Instances}\label{sec:instances}

This section is entirely combinatorial: no algebra is used beyond the fact that
each domain carries a $w$ and each constraint relation is a subuniverse of a
product of domains. It is also where the most representation decisions are
forced, because instances are syntactic objects manipulated by operations ---
weakening, covering, renaming --- that are usually described in words.

\subsection{Instances and solutions}

\begin{definition}[Instance]\label{def:instance}
Fix a finite set $V$ of variables and, for each $x\in V$, an algebra
$\mathbf D_{x}\in\Vn$. A \emph{constraint} is a pair
$C=(\bar x,R)$ where $\bar x=(x_{1},\dots,x_{m})$ is a tuple of variables and
$R\le\mathbf D_{x_{1}}\times\dots\times\mathbf D_{x_{m}}$. An
\emph{instance} $\mathcal I$ is a finite list of constraints; we write
$C\in\mathcal I$, and $\Var(C)$, $\Var(\mathcal I)$ for the variables
occurring. A \emph{subinstance} is a sublist.

A \emph{solution} is a map $s$ with $s(x)\in D_{x}$ for all $x$ such that
$(s(x_{1}),\dots,s(x_{m}))\in R$ for every constraint. The solution set of
$\mathcal I$ is \emph{subdirect} if for every $x$ and every $a\in D_{x}$ there
is a solution with $s(x)=a$.
\end{definition}

\begin{hazard}[Lists, not sets; tuples, not sets of variables]\label{haz:instance-list}
Two representation choices are forced by the way the arguments below go, and
both differ from the naive one.

\emph{Constraints form a list, not a set.} Proofs delete ``a constraint
$C\in\mathcal I$'' and reason about $\mathcal I\setminus\{C\}$, and the same
relation may occur twice on different variable tuples --- and, in the expanded
coverings of \S\ref{sec:coverings}, the same relation on the \emph{same}
tuple can legitimately occur twice. Multiset or list semantics is required.

\emph{A constraint's scope is a tuple, not a set.} Repeated variables in a
scope are allowed and occur --- $\zeta(x,x',z)$ in
\S\ref{sec:main-statements} arises by splitting one variable into two, and
$R(x,x,\dots,x)$ appears explicitly in Definition~\ref{def:expanded}. So
$\Var(C)$ is the image of the scope, and $|{\Var(C)}|$ may be smaller than the
arity.

Once repeated variables are allowed, every semantic operation below that names
a \emph{variable} rather than a \emph{position} --- $\proj_{z}(C)$ in
$1$-consistency, $\proj_{z_{i},z_{i+1}}(C)$ along a path, $\ConOne(C,x)$,
adjacency at a common variable, weakening by a set of variable names, and the
projections in Definition~\ref{def:instance-linked}(ii) --- must be read as
pulled back along the equal-variable diagonal first. Otherwise a constraint
$R(x,x)$ has a full first-coordinate projection even when $R$ contains no
diagonal tuple and the constraint admits no value for $x$. Fixing this notation
uniformly is work this document has not yet done; see \S\ref{sec:status}.
\end{hazard}

\begin{definition}[Relations defined by an instance]\label{def:instance-relation}
For $x_{1},\dots,x_{m}\in\Var(\mathcal I)$, $\mathcal I(x_{1},\dots,x_{m})$ is
the set of $(a_{1},\dots,a_{m})$ such that $\mathcal I$ has a solution with
$s(x_{i})=a_{i}$ for all $i$. It is a subuniverse of
$\mathbf D_{x_{1}}\times\dots\times\mathbf D_{x_{m}}$, being defined by a
primitive positive formula over the constraint relations.
\end{definition}

\subsection{Reductions}

\begin{definition}[Reduction]\label{def:reduction}
A \emph{reduction} for $\mathcal I$ is a map $D^{(\top)}$ assigning to each
$x\in\Var(\mathcal I)$ a subuniverse $D^{(\top)}_{x}\subseteq D_{x}$, possibly
empty. It is \emph{nonempty} if every $D^{(\top)}_{x}\neq\varnothing$. The
instance $\mathcal I^{(\top)}$ is $\mathcal I$ with each domain replaced by
$D^{(\top)}_{x}$ and each constraint relation intersected with the
corresponding box. For reductions we write $D^{(\bot)}\lll D^{(\top)}$ and
$D^{(\bot)}\subte{T}D^{(\top)}$ when the corresponding relation holds at every
variable.
\end{definition}

\begin{hazard}[Empty reductions]\label{haz:empty-reduction}
Reductions are allowed to be empty at some or all variables --- the maximal
$1$-consistent reduction below $D^{(B,\top)}$ in the proof of
Theorem~\ref{thm:main-inductive} is constructed exactly so that it may come out
empty, and the case split on whether it does is the branch point of the whole
argument. But $D^{(\bot)}\lll D^{(\top)}$ requires nonemptiness at every
variable (Convention~\ref{conv:empty}). Both notions are needed and they must
be different objects.
\end{hazard}

\subsection{Consistency}

\begin{definition}[Paths]\label{def:path}
A \emph{path} in $\mathcal I$ is an alternating sequence
$z_{1}-C_{1}-z_{2}-\dots-C_{l-1}-z_{l}$ with $z_{i},z_{i+1}\in\Var(C_{i})$.
It \emph{connects $b$ and $c$} if there are $a_{i}\in D_{z_{i}}$ with
$a_{1}=b$, $a_{l}=c$, and $(a_{i},a_{i+1})\in\proj_{z_{i},z_{i+1}}(C_{i})$ for
every $i$.
$\mathcal I$ is a \emph{tree-instance} if it has no path with $l\ge 3$,
$z_{1}=z_{l}$ and $C_{1},\dots,C_{l-1}$ pairwise distinct.
\end{definition}

\begin{definition}[Consistency]\label{def:consistency}
$\mathcal I$ is \emph{$1$-consistent} if $\proj_{z}(C)=D_{z}$ for every
constraint $C$ and every $z\in\Var(C)$. A reduction $D^{(\top)}$ is
$1$-consistent for $\mathcal I$ if $\mathcal I^{(\top)}$ is $1$-consistent.
$\mathcal I$ is \emph{cycle-consistent} if it is $1$-consistent and every path
from $z$ to $z$ connects $a$ to $a$, for every variable $z$ and every
$a\in D_{z}$.
\end{definition}

\begin{definition}[Linked, fragmented, irreducible]\label{def:instance-linked}
$\mathcal I$ is \emph{linked} if for every $z\in\Var(\mathcal I)$ and every
$a,b\in D_{z}$ some path from $z$ to $z$ connects $a$ and $b$.
$\mathcal I$ is \emph{fragmented} if $\Var(\mathcal I)$ splits into two disjoint
nonempty sets $X_{1},X_{2}$ with $\Var(C)\subseteq X_{1}$ or
$\Var(C)\subseteq X_{2}$ for every $C$.
$\mathcal I$ is \emph{irreducible} if there is no instance $\mathcal I'$ with
\begin{enumerate}\romanlabels\tightlist
\item $\Var(\mathcal I')\subseteq\Var(\mathcal I)$;
\item every constraint of $\mathcal I'$ is a projection of a constraint of
  $\mathcal I$ onto some of its variables;
\item $\mathcal I'$ is not fragmented;
\item $\mathcal I'$ is not linked;
\item the solution set of $\mathcal I'$ is not subdirect.
\end{enumerate}
\end{definition}

\begin{hazard}[Irreducibility is a $\Pi_{1}$ condition over a large family]
\label{haz:irreducible-instance}
Conditions (i)--(v) quantify over \emph{all} instances built from projections
of constraints of $\mathcal I$ --- a finite but exponentially large family
(any multiset of projections of any constraints onto any subsets of their
variables). It is finite, so the predicate is decidable, but it is not
something to unfold. Every use in \S\ref{sec:main-statements} goes through
Lemma~\ref{lem:crucial-irreducible} and \emph{never} through the definition.
\end{hazard}

\subsection{Weakening and cruciality}

\begin{definition}[Weakening]\label{def:weakening}
A constraint $R_{1}(y_{1},\dots,y_{t})$ is \emph{weaker or equivalent to}
$R_{2}(z_{1},\dots,z_{s})$ if $\{y_{i}\}\subseteq\{z_{j}\}$ and
$R_{2}(\bar z)$ implies $R_{1}(\bar y)$; it is \emph{weaker} if in addition
$R_{1}(\bar y)$ does not imply $R_{2}(\bar z)$. \emph{Weakening} a constraint
$C$ in $\mathcal I$ means replacing $C$ by all constraints weaker than it. An
instance $\mathcal I'$ is a \emph{weakening} of $\mathcal I$ if
$\Var(\mathcal I')\subseteq\Var(\mathcal I)$ and every constraint of
$\mathcal I'$ is weaker or equivalent to a constraint of $\mathcal I$.
\end{definition}

\begin{hazard}[``Replace by all weaker constraints'', and its termination]
\label{haz:weakening}
The operation is finite but enormous: there are $2^{|D|^{t}}$ relations weaker
than a given one on a given scope, and one must also range over all sub-scopes.
No algorithm performs it --- it appears only inside proofs, as a device for
reaching a crucial instance (Remark~\ref{rem:get-crucial}). Weakening is
therefore defined above as a \emph{relation between instances}, not as a
function producing one.

The termination of Remark~\ref{rem:get-crucial} is a genuine obligation and is
\emph{not} discharged here. Two things make it delicate, and both must be
faced before ``we cannot weaken forever'' --- which is used three times in
\S\ref{sec:main-statements} --- can be relied on.

\emph{The obvious measure points the wrong way.} A properly weaker constraint
admits \emph{more} tuples, so the total number of tuples in all constraint
relations \emph{increases}. A measure with the right orientation would be a
rank in the finite implication order on constraints, or the number of
assignments forbidden by the instance.

\emph{Replacing a constraint by \emph{all} its strict weakenings need not make
progress.} The conjunction of all strict weakenings of a relation can be the
relation itself --- intersecting all one-tuple extensions of $R$ recovers $R$.
So the construction must choose one proper weakening, or a finite antichain
with a proved rank decrease, and the choice must be part of the statement.

See \S\ref{sec:status}.
\end{hazard}

\begin{definition}[Crucial]\label{def:crucial}
Let $D^{(\top)}$ be a reduction for $\mathcal I$. A variable $y_{i}$ of a
constraint $R(y_{1},\dots,y_{t})$ is \emph{dummy} if $R$ does not depend on its
$i$-th coordinate. A constraint $C\in\mathcal I$ is \emph{crucial in
$D^{(\top)}$} if it has no dummy variables, $\mathcal I^{(\top)}$ has no
solution, and weakening $C$ in $\mathcal I$ yields an instance with a solution
in $D^{(\top)}$. The instance $\mathcal I$ is \emph{crucial in $D^{(\top)}$} if
it has at least one constraint and all of its constraints are crucial in
$D^{(\top)}$.
\end{definition}

\begin{remark}[Reaching a crucial instance]\label{rem:get-crucial}
If $\mathcal I^{(\top)}$ has no solution, then iteratively weakening
constraints that are not crucial, dropping dummy variables as they appear,
terminates in an instance crucial in $D^{(\top)}$. Every relation so introduced
is still a subuniverse of the corresponding product of domains. The termination
is Caveat~\ref{haz:weakening} and is not proved here.
\end{remark}

\subsection{Coverings}\label{sec:coverings}

\begin{definition}[Expanded covering]\label{def:expanded}
$\mathcal I'$ is an \emph{expanded covering} of $\mathcal I$, written
$\mathcal I'\in\Expanded(\mathcal I)$, if there is
$S\colon\Var(\mathcal I')\to\Var(\mathcal I)$ with
\begin{enumerate}\romanlabels\tightlist
\item $S(x)=x$ for $x\in\Var(\mathcal I)\cap\Var(\mathcal I')$;
\item $D_{x}=D_{S(x)}$ for every $x\in\Var(\mathcal I')$;
\item for every constraint $R(x_{1},\dots,x_{m})$ of $\mathcal I'$, either
  $S(x_{1}),\dots,S(x_{m})$ are pairwise distinct and
  $R(S(x_{1}),\dots,S(x_{m}))$ is weaker or equivalent to a constraint of
  $\mathcal I$, or $S(x_{1})=\dots=S(x_{m})$ and
  $\{(a,\dots,a):a\in D_{x_{1}}\}\subseteq R$.
\end{enumerate}
It is a \emph{covering} if moreover $R(S(x_{1}),\dots,S(x_{m}))\in\mathcal I$
for every constraint; a \emph{tree-covering} if it is a covering and a
tree-instance. $S(x)$ is the \emph{parent} of $x$, and $x$ a \emph{child}.
\end{definition}

\begin{proposition}[Basic properties]\label{prop:covering-basics}
Let $\mathcal I'$ be an expanded covering of $\mathcal I$.
\begin{enumerate}\alphlabels\tightlist
\item Replacing every $x$ by $S(x)$ and deleting the constraints
  $R(x,\dots,x)$ yields a weakening of $\mathcal I$.
\item A weakening is an expanded covering with $S=\mathrm{id}$.
\item Every solution of $\mathcal I$ lifts to a solution of $\mathcal I'$.
\item If $\mathcal I$ is $1$-consistent and $\mathcal I'$ is a tree-covering,
  the solution set of $\mathcal I'$ is subdirect.
\item The union of two expanded coverings is an expanded covering.
\item An expanded covering of an expanded covering is an expanded covering.
\item An expanded covering of a cycle-consistent irreducible instance is
  cycle-consistent and irreducible.
\item Every reduction of $\mathcal I$ extends to $\mathcal I'$, and a
  $1$-consistent one extends to a $1$-consistent one.
\end{enumerate}
\end{proposition}

\begin{hazard}[Variable renaming carries obligations]\label{haz:renaming}
The constructions of \S\ref{sec:main-statements} say ``rename the variables so
that the only common variable of $\Upsilon_{x}$ and $\mathcal J$ is $x$'', and
form conjunctions $\bigwedge_{\mathcal J\in\Omega}\mathcal J$ of instances
whose variable sets must first be made disjoint. Renaming is not free: each of
$1$-consistency, cycle-consistency, cruciality and being a covering has to be
shown invariant under it. Those four facts are used throughout
\S\ref{sec:main-statements} and are nowhere stated.
\end{hazard}

\subsection{Induced congruences and connectedness}

\begin{definition}[$\ConOne$]\label{def:conone}
For $R$ of arity $m$ and $i\in[m]$, $\ConOne(R,i)$ is the binary relation
\[
  \{(y,y')\;:\;\exists\,\bar x\ \ R(\bar x[i\mapsto y])\wedge R(\bar x[i\mapsto y'])\}.
\]
For a constraint $C=R(x_{1},\dots,x_{m})$ we write $\ConOne(C,x_{i})$ for
$\ConOne(R,i)$. For an instance, $\Con(\mathcal I,x)=\{\ConOne(C,x):C\in
\mathcal I\}$ and $\Con(\mathcal I)=\bigcup_{x}\Con(\mathcal I,x)$.
\end{definition}

\begin{lemma}\label{lem:conone-basics}
The $i$-th variable of $R$ is rectangular if and only if $R$ is stable under
$\ConOne(R,i)$. If $R$ is subdirect and its $i$-th variable is rectangular,
then $\ConOne(R,i)$ is a congruence.
\end{lemma}

\begin{hazard}\label{haz:conone-not-congruence}
$\ConOne(R,i)$ is always reflexive and symmetric but is \emph{not} transitive
in general, hence not always a congruence; Lemma~\ref{lem:conone-basics} is
what makes it one, and its hypotheses (subdirect \emph{and} rectangular at $i$)
are both needed. The notation invites the reader to assume congruence-hood, and
every use below must check the hypotheses. Note also that
$\ConOne(C,x_{i})$ depends on the \emph{position} $i$, not on the variable, so
it is ill-defined if $x$ occurs twice in the scope
(Caveat~\ref{haz:instance-list}).
\end{hazard}

\begin{definition}[Types of relations and instances]\label{def:rel-type}
$R$ is \emph{of PC type} (resp.\ \emph{linear type}) if $R$ is rectangular and
every $\ConOne(R,i)$ is a PC (resp.\ linear) congruence. An instance is of PC
or linear type if all its constraints are.
\end{definition}

\begin{definition}[Adjacency, connectedness]\label{def:connected}
Two congruences $\sigma_{1},\sigma_{2}$ on $\mathbf D_{x}$ are \emph{adjacent}
if there is a reflexive bridge from $\sigma_{1}$ to $\sigma_{2}$. Two
rectangular constraints $C_{1},C_{2}$ are \emph{adjacent in a common variable
$x$} if $\ConOne(C_{1},x)$ and $\ConOne(C_{2},x)$ are adjacent. An instance
$\mathcal I$ is \emph{connected} if all its constraints are rectangular, all
congruences in $\Con(\mathcal I)$ are irreducible, and the graph on the
constraints with edges the adjacent pairs is connected.
\end{definition}

Note that every proper congruence is adjacent to itself, via
$\delta(x_{1},x_{2},x_{3},x_{4})=\sigma(x_{1},x_{3})\wedge\sigma(x_{2},x_{4})$.
